% Layout PI – tema roxo/dourado (inspirado em Pi Network / guias de mineração)
\usepackage{xcolor}
\usepackage{titlesec}
\usepackage{fancyhdr}
\usepackage{pifont}

% Paleta: fundo roxo escuro, texto claro, destaque dourado
\definecolor{pibg}{RGB}{35,20,55}
\definecolor{pitext}{RGB}{255,252,245}
\definecolor{gold}{RGB}{255,193,7}
\AtBeginDocument{%
  \pagecolor{pibg}%
  \color{pitext}%
}

% Títulos em dourado (estilo capa Pi)
\titleformat{\section}{\color{gold}\Large\bfseries}{\thesection}{1em}{}
\titleformat{\subsection}{\color{gold!95}\large\bfseries}{\thesubsection}{1em}{}
\titleformat{\subsubsection}{\color{gold!90}\normalsize\bfseries}{\thesubsubsection}{1em}{}

% Listas com marcadores: check (✓) em dourado; numeradas: número em dourado
\renewcommand{\labelitemi}{\textcolor{gold}{\ding{51}}}
\renewcommand{\labelenumi}{\textcolor{gold}{\theenumi.}}

% Rodapé em dourado (estilo faixa da capa Pi)
\pagestyle{fancy}
\fancyhf{}
\fancyfoot[C]{\color{gold}\small Creative Engine\quad$\cdot$\quad\thepage}
\renewcommand{\headrulewidth}{0pt}
\renewcommand{\footrulewidth}{0pt}

% Links em dourado
\AtBeginDocument{\hypersetup{colorlinks=true,linkcolor=gold!90,urlcolor=gold!80,citecolor=gold!80}}
